\section{Finite reference frames and verdict accounting}
\label{app:accounting}

Let $\mathcal J_i$ be the nonempty set of criterion slots for evaluation unit
$i$, with $k_i=|\mathcal J_i|$, $N$ units, and
$M=\sum_i k_i$ slots. Slots belong to exactly one unit. The released reference
and primary prediction for slot $u$ are $Y_u,P_u\in\{0,1\}$, where one denotes
PASS. Throughout this appendix they are fixed finite-population values. Define
\begin{equation}
 Z_i=\prod_{u\in\mathcal J_i}Y_u,
 \qquad \widehat Z_i^0=\prod_{u\in\mathcal J_i}P_u,
 \qquad \theta=\frac{1}{N}\sum_i Z_i.
 \label{app:eq:targets}
\end{equation}
Criterion microagreement is the distinct target
$a=M^{-1}\sum_u\mathbf 1\{Y_u=P_u\}$.
An error in the following equations means disagreement with this particular
reference. No equation identifies which of two disagreeing human raters is
substantively correct.

For JudgmentBench full strict, a positive binary item passes when its original
value is one, and a negative-weight occurrence-count item passes when its
original count is zero. The loader checks that their conjunction equals
attaining the rubric's maximum score. Binary-only retains positive binary
items and removes occurrence-count items. This changes both the success event
and the criterion population. The original unit is an output-by-rater
annotation; the alternative-reference analysis instead weights distinct
outputs equally. RuVerBench units use the criteria retained in its release.

\subsection{Certificates and update rules}
\label{app:certificates}

Let $Q\subseteq\bigcup_i\mathcal J_i$ be the queried slots and
$Q_i=Q\cap\mathcal J_i$. In the query-only information model, a FAIL
certificate is an observed reference failure, and a PASS certificate requires
observing every slot as PASS. Write
\begin{align}
 C_i^-&=\mathbf1\{\exists u\in Q_i:Y_u=0\},\\
 C_i^+&=\mathbf1\{Q_i=\mathcal J_i,\ Z_i=1\},
 \qquad C_i=C_i^-+C_i^+.
 \label{app:eq:certificates}
\end{align}
Eager updating substitutes every queried reference value immediately:
\begin{equation}
 \widehat Z_i^{E}(Q)
 =\prod_{u\in\mathcal J_i}
       \bigl(\mathbf1\{u\in Q\}Y_u+\mathbf1\{u\notin Q\}P_u\bigr).
 \label{app:eq:eager}
\end{equation}
Certificate-gated updating uses $Z_i$ once certified and otherwise retains
$\widehat Z_i^0$:
\begin{equation}
 \widehat Z_i^{G}(Q)=C_iZ_i+(1-C_i)\widehat Z_i^0.
 \label{app:eq:gated}
\end{equation}
Writing $Z_i$ here is shorthand for the value logically determined by a
certificate, not access to an unqueried reference. The implementation obtains
that value from the queried bits. Gating can retain an uncertified initial
PASS; it is not a protocol that publishes only certified outcomes.

\subsection{Acquisition implementations}
\label{app:acquisition}

The seven acquisition rules use a common slot map. Uniform criterion
sampling follows a global random priority. Three serial short-circuit rules
visit tasks in the stated order and visit their criteria in natural slot
order, random order, or predicted-FAIL-first order, respectively. They stop
the current task at its first observed reference failure or after completing
its rubric. Disagreement-only visits flagged criteria within each task and
stops after the flagged pool is exhausted; it can leave budget unused.
Disagreement-then-random uses that same flagged prefix, followed by a global
random ordering of all unflagged criteria. Disagreement-prioritized short
circuit visits, within a task, flagged slots first, then unflagged predicted
FAIL slots, then the remaining predicted PASS slots; its stopping rule is
the same serial certification rule.

For each DR primary, the auxiliaries are the other two saved DR judges
among Gemini 3.1 Pro, Qwen-plus and GPT-5.4 low. Each AC primary has the
other saved AC judge as its auxiliary: Qwen-plus or DeepSeek v4-pro.
JudgmentBench pairs GPT-5.4 with GPT-5.4-mini within the same scoring target.
The flag is the union of disagreements with available auxiliaries. Their
predictions are treated as available signals; their acquisition cost is not
included in the reference-bit budget.

The three task orders are release-relative order, ascending primary
predicted-failure count, and a random task order. Seeds 1--31 randomize ties
within the declared priorities. A separate canonical seed-zero realization
uses natural slot order for priority ties while retaining random priorities
for explicitly random policies and the random task order. Purpose-specific
streams derive from root seed 20260913, the replicate seed and a SHA256 hash
of the stream's name. Release-relative order refers to the loader's retained
unit ordering, not a randomized sample of deployment tasks. The analytical
mechanism calculations below hold this order fixed.

\subsection{Exact event decomposition}
\label{app:events}

Define predicted-failure slots $F_i=\{u\in\mathcal J_i:P_u=0\}$ and
reference-failure slots $T_i=\{u\in\mathcal J_i:Y_u=0\}$. An initially
correct FAIL is susceptible to a new false PASS only if $F_i$ and $T_i$ are
nonempty and disjoint. Its new-error event is
\begin{equation}
 A_i(Q)=\mathbf1\{F_i\ne\varnothing,\ T_i\ne\varnothing,
 F_i\cap T_i=\varnothing,\ F_i\subseteq Q,
 T_i\cap Q=\varnothing\}.
 \label{app:eq:introduced}
\end{equation}
An initially false FAIL has $T_i=\varnothing$ and
$F_i\ne\varnothing$. Eager updating repairs it once all its predicted
failures have been checked. Gating delays that repair until every criterion
has been checked. The delayed-repair event is therefore
\begin{equation}
 D_i(Q)=\mathbf1\{T_i=\varnothing,\ F_i\ne\varnothing,
 F_i\subseteq Q,\ \mathcal J_i\not\subseteq Q\}.
 \label{app:eq:delayed}
\end{equation}
Let $A(Q)=\sum_i A_i(Q)$, $D(Q)=\sum_i D_i(Q)$, and
$E_r(Q)=\sum_i\mathbf1\{\widehat Z_i^r(Q)\ne Z_i\}$ for
$r\in\{E,G\}$. Then
\begin{equation}
 E_G(Q)-E_E(Q)=D(Q)-A(Q).
 \label{app:eq:update-identity}
\end{equation}

\paragraph{Derivation.}
Initially correct PASS units have every prediction and reference equal to
one, so either update remains correct. An initially false PASS is repaired
exactly when a reference failure is queried, which also certifies FAIL; both
updates therefore agree. For an initially correct FAIL, an unqueried
predicted failure or a queried reference failure prevents eager PASS. Eager
can introduce a false PASS precisely in \eqref{app:eq:introduced}, while
gating retains the correct FAIL. For an initially false FAIL, the two rules
disagree precisely in \eqref{app:eq:delayed}. Summing these four exhaustive
initial classes proves \eqref{app:eq:update-identity}. In particular, correct
reference substitutions cannot introduce a false FAIL under conjunction.

If false PASS and false FAIL have fixed nonnegative costs $c_{\rm FP}$ and
$c_{\rm FF}$, their endpoint loss difference is
$c_{\rm FF}D-c_{\rm FP}A$. Thus equal-weight error comparisons are one
particular loss choice. Along a growing query transcript, gated error cannot
increase: every changed verdict is certified relative to the fixed
reference. Its difference from eager error can increase because eager may
repair false FAILs earlier. At zero queries and at a census the difference
is zero.

Both rules agree on untouched and certified units. If $U(Q)$ units have been
partially queried but remain uncertified, then
$|E_G(Q)-E_E(Q)|\leq U(Q)$. A serial short-circuit policy that completes or
certifies a task before starting the next has $U(Q)\leq1$ at every prefix.
This explains why changing the update rule has little room to change the
endpoint under that acquisition structure, even when dispersed querying has
a large effect.

\section{Finite-budget expectations over query order}
\label{app:expectations}

For uniform sampling of $h$ slots without replacement from a group of size
$G$, the probability of including $r$ specified slots and excluding $f$
other specified slots is
\begin{equation}
 H(G,h;r,f)=\frac{\binom{G-r-f}{h-r}}{\binom{G}{h}}.
 \label{app:eq:hypergeom-event}
\end{equation}
The specified inclusion and exclusion sets are disjoint. We set the
probability to zero when $r+f>G$, $r>h$, or $f>G-h$; in the empty case
$H(0,0;0,0)=1$. Equation~\eqref{app:eq:hypergeom-event} follows by choosing
the other $h-r$ observed slots from the $G-r-f$ unrestricted slots. It is a
sampling calculation and assumes no independence of reference errors.

For a uniform $m$-of-$M$ criterion query set, the two expectations are
\begin{align}
 \mathbb E[A(Q)]
 &=\sum_{i:\,F_i,T_i\ne\varnothing,\,F_i\cap T_i=\varnothing}
     H(M,m;|F_i|,|T_i|),\\
 \mathbb E[D(Q)]
 &=\sum_{i:\,T_i=\varnothing,\,F_i\ne\varnothing}
     \{H(M,m;|F_i|,0)-H(M,m;k_i,0)\}.
 \label{app:eq:uniform-expectations}
\end{align}
The complete-rubric event is contained in the event of observing every
predicted failure, which justifies the subtraction in the second line.
Linearity of expectation gives the expected update difference even though
events for different tasks are dependent under a fixed global budget.

\subsection{Prioritized groups}
\label{app:priority-groups}

The same calculation applies to a predetermined sequence of disjoint query
groups $\mathcal G_1,\ldots,\mathcal G_L$. Every slot in an earlier group is
queried before any slot in a later group, and order within each group is
uniformly random. At an interior cap $m$, let $\ell$ be the group satisfying
\[
 \sum_{g<\ell}|\mathcal G_g|\leq m
 <\sum_{g\leq\ell}|\mathcal G_g|,
 \qquad h=m-\sum_{g<\ell}|\mathcal G_g|.
\]
For a required set $R$ and forbidden set $S$, the event
$R\subseteq Q$, $S\cap Q=\varnothing$ has probability zero if a required
slot lies after $\ell$, a forbidden slot lies before $\ell$, or the two sets
overlap. Otherwise its probability is
\begin{equation}
 H\bigl(|\mathcal G_\ell|,h;
         |R\cap\mathcal G_\ell|,|S\cap\mathcal G_\ell|\bigr).
 \label{app:eq:group-event}
\end{equation}
At a complete group boundary, the next group has $h=0$. At a census the
event is determined directly by the complete query set. Substituting this
event probability for each term in
\eqref{app:eq:uniform-expectations} gives the grouped-order expectations.

For disagreement-then-random, a criterion is flagged when any available
auxiliary judge disagrees with the primary prediction. Each task's flagged
criteria form one group, tasks follow release order, and the final group
contains all unflagged criteria. The flagged tier is therefore not a single
global random sample. These groups depend on saved predictions and task
order, not hidden references. The analytical grid covers uniform criterion
and disagreement-then-random querying at 101 shares from zero to one. It
does not apply the predetermined-group formula to adaptive short-circuit
paths. Budget caps use the implementation's nearest-integer rule
$m=\operatorname{round}(bM)$, with ties to an even integer.

All expectations condition on the full reference, predictions and priority
groups. They explain the observed frames rather than provide a free
prospective policy selector. The 31 random-order means provide a Monte Carlo
check against these expectations; their standard error is the sample
standard deviation divided by $\sqrt{31}$, not uncertainty from sampling new
tasks.

\subsection{Matched-margin alternatives}
\label{app:permutations}

We construct 64 alternative assignments in each of nine analysis cells by
permuting reference labels within strata defined by task length, primary
prediction and criterion category or scoring mode. This preserves every
stratum's prediction/reference confusion counts, the task slots and lengths,
and all primary and auxiliary predictions. Root seed 20261007 and a
dataset-specific hash define the random streams. Both analytical policies
are evaluated at shares $0.05,0.20,0.50,0.84,0.95$, including originally
observed focal and peak regions. An alternative assignment with a different
update effect demonstrates that the retained summaries do not determine
that effect. These are retrospective counterexamples, not a null test for
random reference noise. The permutations can also change reference task
pass counts, initial task errors and their joint configuration, so they do
not isolate a causal effect of error correlation.

\section{Certification and repair under different allocations}
\label{app:allocation}

Let $E_0=\sum_i\mathbf1\{\widehat Z_i^0\ne Z_i\}$. For a query set $Q$,
write $C_W(Q)$ and $C_R(Q)$ for the numbers of certified units that were
initially wrong and initially right. Then
\begin{align}
 C(Q)&=C_W(Q)+C_R(Q),\\
 E_G(Q)&=E_0-C_W(Q),\\
 E_E(Q)&=E_0-C_W(Q)-D(Q)+A(Q).
 \label{app:eq:allocation-identity}
\end{align}
Every certified initially wrong unit is repaired under either rule. The
only further eager repairs are the uncertified false-FAIL repairs counted
by $D$, and its only new errors are those counted by $A$. This proves the
last two identities directly from the initial classes.

For two allocations applied to the same frame, define every signed
difference as mixed allocation minus uniform criterion sampling. Pairing
within a query-order replicate yields
\begin{equation}
 \Delta C=\Delta C_W+\Delta C_R,
 \qquad
 \Delta E_E=-\Delta C_W-\Delta D+\Delta A,
 \qquad
 \Delta E_G=-\Delta C_W.
 \label{app:eq:paired-allocation}
\end{equation}
We calculate these differences before averaging. Their means retain the
additive identities, whereas differences of marginal medians generally do
not. To describe where additional certificates go, we also report the
composition of the \emph{mixed-only} certified set, separately from the
uniform-only set. This gross composition remains meaningful when net
certification gain is zero or negative. More certificates on initially
correct units can provide useful confirmation without repairing an error;
it is not by itself evidence that an allocation is harmful.

The fixed mixed design spends up to $\lfloor m/2\rfloor$ labels on
judge-first serial short circuit, then draws a fresh uniform sample without
replacement from the remaining criterion slots. If short circuit completes
before its phase cap, the unused budget transfers to the random phase. All
revealed labels count toward the common actual total $m$. This contrasts
acquisition while the matched-query comparison in
Appendix~\ref{app:events} contrasts updating. Neither design uses hidden
initial-error classes to choose queries.

\section{Finite-population estimation and intervals}
\label{app:inference}

\subsection{An adaptive first phase followed by fresh sampling}

Set $a_u=\mathbf1\{P_u=Y_u\}$. Let the first-phase transcript reveal a set
$C$ of $c$ criterion labels with known correct count $t_C=\sum_{u\in C}a_u$.
Condition on that complete transcript. The remaining $R=M-c$ slots form a
fixed population. Draw an independent simple random sample $S$ of size $n$
without replacement from it, and observe $x=\sum_{u\in S}a_u$. For $n>0$,
\begin{equation}
 \widehat a=\frac{t_C+R x/n}{M},
 \qquad \mathbb E[\widehat a\mid C,\hbox{first-phase observations}]=a.
 \label{app:eq:conditional-estimator}
\end{equation}
The expectation follows from inclusion probability $n/R$ for every remaining
slot. Conditioning on the observed history accommodates the adaptive first
phase: the second-phase estimate expands only the unobserved remainder.
When $R=0$, the reference agreement is known exactly. When $R>0$ but $n=0$,
we report logical bounds $[t_C/M,(t_C+R)/M]$ and no design-unbiased point
estimate from this formula.

Let $S_R^2$ be the finite-population variance of the remainder's agreement
bits, with denominator $R-1$. For $R>1$ and $n>0$,
\begin{equation}
 \operatorname{Var}(\widehat a\mid\hbox{first phase})
 =\frac{R^2}{M^2}\left(1-\frac{n}{R}\right)\frac{S_R^2}{n}.
 \label{app:eq:fpc}
\end{equation}
The factor $1-n/R$ is the finite-population correction. Census remainders
have zero variance. Where a normal approximation is reported, $n\geq2$
permits the sample variance
$s^2=x(n-x)/[n(n-1)]$ to replace $S_R^2$. That approximation is evaluated
separately from the guaranteed intervals below.

\subsection{Inclusive hypergeometric inversion}
\label{app:hypergeom-ci}

For unknown remainder success count $K$, the sample count has probability
\begin{equation}
 \Pr_K(X=j)=\frac{\binom{K}{j}\binom{R-K}{n-j}}{\binom{R}{n}}.
 \label{app:eq:hg-pmf}
\end{equation}
Retain integer counts consistent with the sample and both inclusive tails:
\begin{equation}
 \mathcal K_\alpha(x)=
 \left\{K\in\{x,\ldots,R-n+x\}:
 \Pr_K(X\geq x)\geq\frac{\alpha}{2},\quad
 \Pr_K(X\leq x)\geq\frac{\alpha}{2}\right\}.
 \label{app:eq:hg-set}
\end{equation}
Let $K_L$ and $K_U$ be its minimum and maximum. The microagreement interval
is $[(t_C+K_L)/M,(t_C+K_U)/M]$. For every fixed true $K$, rejection by either
tail occurs with probability at most $\alpha/2$; a union bound gives coverage
at least $1-\alpha$. Discreteness can make it conservative. This is a
conditional, fixed-budget randomization guarantee with the reference held
fixed, not simultaneous coverage of every budget or validity under optional
stopping. At $n=0$ the count range is $[0,R]$, and at $n=R$ it is $[x,x]$.

The implementation compares ordinary tails numerically, but recomputes
near-threshold decisions using integer binomial coefficients and rational
$\alpha/2$. A tolerance triggers exact calculation rather than accepting an
endpoint. For example, with $R=16,n=2,x=0,\alpha=0.05$,
$K=13$ satisfies $\Pr(X=0)=\binom{3}{2}/\binom{16}{2}=0.025$ exactly and
must remain in the interval $[0,13]$. This resolves the floating-point
boundary case without altering the mathematical acceptance rule.

\subsection{Prediction strata and a charged pilot}
\label{app:strata}

Prediction strata are the nonempty sets with $P_u=0$ or $P_u=1$, known before
reference acquisition. Proportional allocation reserves one draw per
nonempty stratum when affordable, distributes remaining draws proportional
to stratum size, respects capacities, and resolves integer remainders by
descending fractional part with prediction-zero first on ties. Otherwise it
falls back to uniform sampling.

For the charged-pilot design, the initial allocation is
$\min(2,N_h)$ labels in stratum $h$. The design reserves at least one later
draw in each stratum not exhausted by that initial allocation. If the cap
cannot fund both stages, it falls back to uniform sampling. Otherwise the
pilot target is the larger of this minimum and $\lfloor0.2m\rfloor$, capped
to preserve the later draws. Additional pilot slots are allocated by stratum
size, leaving an unobserved slot in each initially noncensused stratum.
After pilot counts $p_h$ and observed agreement totals $t_h$, define
\[
 \widetilde a_h=\frac{t_h+1/2}{p_h+1},\qquad
 R_h=N_h-p_h,\qquad
 w_h=R_h\sqrt{\widetilde a_h(1-\widetilde a_h)}.
\]
The remaining budget is allocated by weights $w_h$, with a minimum of one
draw from each nonempty remainder and its capacity as an upper bound. The
same capacity-constrained proportional allocation and largest-remainder
rule determine integer counts. Pilot and remainder samples use separate
random-priority streams. All pilot labels are charged to the budget.

Conditional on the pilot, let the remainder sample in stratum $h$ have size
$n_h$ and agreement total $x_h$. Then
\begin{equation}
 \widehat a_{\rm strat}
 =\frac{\sum_h t_h+\sum_{h:R_h>0}R_hx_h/n_h}{M}
 \label{app:eq:strat-estimator}
\end{equation}
is conditionally unbiased. If $H$ strata are not censused, we obtain a count
interval in each at error level $\alpha/H$, sum their lower and upper bounds
and add the known pilot total. The union bound proves conditional coverage
at least $1-\alpha$ for the overall mean; independence between stratum
interval events is not required. Censused strata contribute their exact
counts. The accompanying normal approximation sums the variances in
\eqref{app:eq:fpc}, stratum by stratum, and clips endpoints to $[0,1]$.
It is unavailable when a noncensused stratum has fewer than two sampled
slots. These design-based controls are not an implementation of StratPPI.

At each of six shares $0.006,0.05,0.20,0.50,0.60,1$, we report bias, empirical
RMSE, coverage and interval width over 31 query-order replicates in each of
nine analysis cells. RMSE is
$[31^{-1}\sum_{s=1}^{31}(\widehat a_s-a)^2]^{1/2}$.
The replications describe the finite frames; coverage guarantees derive
from the sampling design and inversion, not the number of observed covered
replicates. Conservative simultaneous stratum bounds can be wider despite
a smaller empirical RMSE.

\subsection{Whole-task sampling}
\label{app:whole-task}

For estimating $\theta$, the separate whole-task panel fixes the number of
sampled evaluation units in advance and draws them uniformly without
replacement. Fully revealing an included unit and short-circuiting after
its first reference failure yield the same certified $Z_i$. Their sample
means and finite-population task-count intervals are consequently identical,
although their criterion-query costs differ. Here sample size is fixed and
label cost is random. This experiment is distinct from stopping task
sampling at a fixed label cap, which would generally alter inclusion
probabilities and require its own inference.

\section{Alternative references and task composition}
\label{app:sensitivity}

\subsection{Matched reference panels}
\label{app:reference}

The matched frame contains 206 outputs with at least two distinct raters.
The corresponding 431 annotations are distinguished from the original
1,539 annotation units. Outputs are sorted by base-task and output identity.
For each output, distinct-rater annotations are sorted by the SHA256 digest
of the literal string
\texttt{reference-panel-v1|output|rater|annotation}; the first two complete
rubric vectors form panels A and B. A separate ordering with salt
\texttt{judge-anchor-v1} chooses the annotation supplying the saved primary
and auxiliary predictions. The selection rule does not inspect label values
or agreement. Rubric-slot identities are checked before alignment, and the
same selections are used for both targets and both judges. A and B thus
refer to symmetric selections across outputs, not globally ranked raters or
a chronological improvement in reference quality.

Crossed comparisons distinguish the reference used for acquisition from
the reference used for reading and evaluating the resulting query set.
For acquisition panel $a\in\{A,B\}$, an adaptive policy obtains $Q^{(a)}$
using that panel's answers. For evaluation panel $e\in\{A,B\}$, queried
slots are re-read using $Y^{(e)}$, both updates are computed from those
answers, and residual disagreement is measured against the complete
$Y^{(e)}$. Off-diagonal comparisons therefore describe counterfactual
re-reading of fixed queried items. They do not retain panel A answers while
silently treating panel B as truth. Nonadaptive policies have identical
query sets for either acquisition panel. Adaptive short circuit can change
its stopping decisions and actual query count.

The update grid includes uniform criterion, disagreement-then-random and
judge-first short circuit at all 101 budget shares and seeds 1--31. The
allocation grid compares the mixed and uniform designs at the six shares
listed above, with equal actual query counts within each contrast. The four
crossed means share outputs and references and are not four independent
replications. Moving from the original frame to the matched frame also
changes composition, weighting and judge anchoring; only within-frame,
fixed-prediction contrasts isolate the reference-panel change.

\subsection{Base-task composition replay}
\label{app:composition}

For the two JudgmentBench GPT-5.4 targets and DR Gemini, we generate 399
bootstrap compositions with root seed 20261008. Each replicate draws the
original number of base tasks with replacement and retains every output,
rater and criterion belonging to a selected task. The two JudgmentBench
targets share base-task multiplicities; DR has a separate random stream.
Original release-relative unit order is preserved by visiting original
units and cloning each according to its base-task multiplicity. Shared
base-task/copy identifiers retain dependence within a replicated task.
Each composition gets its own $N$, $M$ and recomputed rounded budget caps.

For updating, every composition uses the exact order expectations in
Appendix~\ref{app:expectations} for both analytical policies over the full
101-point grid. For allocation, each composition reruns the 20\% mixed versus
uniform contrast using the same fixed 31-seed panel, taking paired
differences before their mean. Thus the latter ranges include sensitivity
of a fixed finite Monte Carlo panel, rather than separately estimating its
Monte Carlo error. Quantiles use linear interpolation. We show pointwise
2.5th and 97.5th percentiles in counts and per 100 evaluation units, and
retain each replicate's recomputed peak and its location. The archived
approximate simultaneous construction uses the 95th percentile of maximum
absolute grid deviations; it is omitted from the main figure because its
constant width obscures the known zero effects at the two endpoints.

This replay evaluates sensitivity to empirical task composition. In
JudgmentBench, 30 curated base tasks do not establish an exchangeable sample
from a new task population. Raters remain fixed, and reference-panel and
composition experiments are conducted separately. Their ranges therefore
do not estimate joint uncertainty over tasks and raters or correct selection
bias. A pointwise range at the original curve's selected peak is also not a
simultaneous guarantee for that curve.

\section{Information disclosure and sharp logical endpoints}
\label{app:disclosure}

\subsection{Query-only identification}

With only queried labels and no restrictions linking unobserved values,
let $n_P=\sum_i C_i^+$ and $n_F=\sum_i C_i^-$. Every uncertified unit has
at least one unobserved slot and no observed failure. Setting all its
unobserved slots to one makes it PASS; setting any one of them to zero makes
it FAIL. Because the unrestricted completions can be chosen independently
by unit, the sharp endpoint range is
\begin{equation}
 \theta\in\left[\frac{n_P}{N},1-\frac{n_F}{N}\right],
 \qquad \text{width}=1-\frac{n_P+n_F}{N}.
 \label{app:eq:query-only}
\end{equation}
Attainable finite-population values have increments $1/N$; the displayed
interval is their endpoint hull. Its width is a logical consequence of
missing information, not a sampling confidence interval. A random sample
can estimate the mean precisely without certifying every population unit.

\subsection{Fixed references with withheld associations}
\label{app:milp}

The disclosure experiment keeps the complete observed labels fixed when
constructing a release and changes what the reconstruction procedure may
see. Visible information includes task slots and lengths, criterion
categories, and the primary predictions. D0 supplies the exact number of
reference PASS bits within each task-length-by-prediction group. D1 refines
those groups by criterion category. D2 additionally restores the actual
reference values attached to all slots of selected whole tasks. These are
reference-containing count summaries and restored links, not additional
label acquisition.

Let $\mathcal B_g$ partition the slots into disclosed groups, with reference
success count $t_g$. Unknown binary assignments $x_u$ and task conjunctions
$z_i$ must satisfy
\begin{align}
 x_u,z_i&\in\{0,1\},\\
 z_i&\leq x_u &&(u\in\mathcal J_i),\\
 \sum_{u\in\mathcal J_i}x_u-z_i&\leq k_i-1 &&(i=1,\ldots,N),\\
 \sum_{u\in\mathcal B_g}x_u&=t_g &&\text{for every disclosed group},\\
 x_u&=Y_u &&\text{for restored slots}.
 \label{app:eq:milp}
\end{align}
The two conjunction inequalities enforce
$z_i=\prod_{u\in\mathcal J_i}x_u$: any zero slot forces $z_i=0$, and all
one slots force $z_i=1$. Minimizing and maximizing $N^{-1}\sum_i z_i$ gives
the sharp finite feasible endpoints. This is a binary mixed-integer
program, not a continuous linear-program relaxation. Every feasible binary
assignment is an allowed completion and vice versa, proving sharpness of
the optima in this information model. Extra disclosures can only remove
completions, and restoring all labels recovers the observed score.

The program separates exactly by task length because every disclosed group
lies within one length stratum. Each endpoint has an attained integer
witness, independently checked equality and conjunction constraints, and a
closed solver objective bound. The solver receives only the constructed
release; hidden labels are used by the publisher to form summaries and
later to check containment and full recovery. The experiment restores
nested whole-task prefixes at fractions $0,0.1,0.25,0.5,0.75,1$, with
$\lfloor fN\rfloor$ tasks, under ten saved random orders per population and
root seed 20260907. The four RuVerBench runs are analyzed separately.
Disclosure-order variation is not sampling uncertainty.

D0 already determines the task-uniform mean-criteria reference score:
\begin{equation}
 \mu=\frac{1}{N}\sum_i\frac{1}{k_i}
                      \sum_{u\in\mathcal J_i}Y_u
     =\frac{1}{N}\sum_g\frac{t_g}{k_g},
 \label{app:eq:linear-control}
\end{equation}
where $k_g$ is the common task length in group $g$. It need not determine
the nonlinear conjunction. Because the predicted task score is visible,
subtracting it from each reference endpoint gives the feasible correction
range. Its sign concerns correction of one judge's score on fixed outputs,
not ranking two solver models. A disclosed reference task-pass total would
also identify the aggregate mean without revealing every criterion link;
the studied disclosure ladder is one sufficient-information comparison,
not a universally minimal release scheme.

\subsection{A complete two-task example}
\label{app:example}

Take two tasks with two binary slots each. With no queries or extra
information, their conjunction mean lies between zero and one. A trusted
total of exactly two PASS bits reduces the feasible means to zero or
$1/2$: assignments $(1,0),(1,0)$ and $(1,1),(0,0)$ attain the two endpoints.
If both first-task slots are then observed as PASS, the trusted total forces
both second-task slots to FAIL, identifying the mean as $1/2$. Without that
total, the same two observations leave a range $[1/2,1]$. Thus a task can
be unqueried yet logically determined by additional information. Conversely,
knowing only that exactly one of the two tasks passes identifies the mean
as $1/2$ without identifying which task passes. These examples distinguish
label acquisition, individual identification and aggregate identification.
The real RuVerBench release contains the task links; the experiment
investigates what its coarsened summaries support, rather than alleging a
missing-link defect in that release.


\clearpage
\section{Additional displays and reproduction}
\label{app:reproduce}
\begin{figure}[!htbp]
\centering
\includegraphics[width=\linewidth,height=0.79\textheight,keepaspectratio]{conjunction_robustness_estimation.pdf}
\caption{Estimation controls at a 20\% criterion budget across all nine cells. Guaranteed interval widths and empirical RMSE are shown separately. The stratified intervals combine finite-population count bounds conservatively; their width alone does not rank point-estimation performance. Pilot labels count toward the same budget. Each empirical quantity summarizes 31 saved orders, while the coverage guarantee follows from the design.}
\label{fig:estimation}
\end{figure}
\FloatBarrier

The released package separates data integrity, scientific replay, and manuscript
validation. Inputs are compact numerical projections with task, output, rater,
criterion and scoring metadata. Original input provenance and content hashes
are in \texttt{data/PROVENANCE.json}; result provenance is stored with the
archived numerical outputs. The following commands run from the repository
root after installing the pinned Python dependencies:
\begin{verbatim}
python -m src.check_package
python -m pytest -q
python -m src.reproduce --mode full --output artifacts/reproduced
python -m src.verify_release --results artifacts/reproduced
python -m src.reproduce_extensions --output artifacts/reproduced-extensions
python -m src.reproduce_extensions --output artifacts/reproduced-extensions --verify-only
python -m src.verify_manuscript_claims --output artifacts/paper-checks/claims.json
python -m src.build_paper
\end{verbatim}
The final command additionally requires Tectonic (tested with version 0.17.0)
and writes a rebuilt PDF separately from the published copy. Initial dependency
installation can require network access; the scientific replay does not call
judge APIs. Rebuilding the PDF can change PDF metadata without changing the
scientific results. The published manuscript includes a claim audit and a
citation audit; the broader technical reading library remains separate from
the selected manuscript bibliography.

Table~\ref{tab:sources} maps the paper's numerical claims to the result groups.
For compactness, paths below are relative to
\texttt{artifacts/expected/}. The executable claim audit records exact row
keys, units, recomputed values and source hashes. It checks selected paper
anchors and accounting identities; the complete scientific reproduction
commands compare the larger archived grids.
\begin{table}[h]
\centering\small
\caption{Numerical source routing for the paper.}
\label{tab:sources}
\begin{tabularx}{\linewidth}{lX}
\toprule
Paper component & Archived source group\\
\midrule
Frames and initial errors & \texttt{conjunction\_policy\_study/study.json} and input tables\\
Query-order curves & \texttt{conjunction\_policy\_study/budget\_curves.csv.gz}\\
Exact expectations and permutations & \texttt{conjunction\_mechanism/}\\
Paired allocation decomposition & \texttt{conjunction\_mechanism/}\\
SRS/mix and whole-task estimates & \texttt{conjunction\_label\_value/}\\
Reference and composition sensitivity & \texttt{conjunction\_robustness/}\\
Stratified estimation controls & \texttt{conjunction\_robustness/}\\
Disclosure endpoints and witnesses & \texttt{information\_disclosure\_replay.json}\\
\bottomrule
\end{tabularx}
\end{table}
